% status: SKELETON
% Chapter 5 of the second edition. Plan: docs/STRUCTURE.md. Rules: AUTHORING.md.
\chapter{The KV Cache: Inference's Central Data Structure}\label{ch:kv-cache}

\begin{ChapterIntent}
Every token a model generates attends to every token before it.
Recomputing their keys and values at each step would make generation
quadratic; caching them makes it linear, and turns the engine into a memory
manager. This chapter derives what is cached and when, computes its size for
real models --- from half a megabyte to a few tens of kilobytes per token,
depending on the architecture --- and shows how the cache, not the weights,
comes to dominate memory capacity at long context and large batch, and the
bytes read per step at long context. Nearly every chapter in
Part~\ref{part:serving} manages this one data structure.
\end{ChapterIntent}

\OpeningSection{At 32K tokens the cache outweighs the model}

\NapkinSection{Napkin math: bytes per token, sequences per GPU}

\section{Why recomputation is quadratic}

\section{What is cached, and when}

\section{State that grows, per request}

\section{Reading the whole cache every step}

\section{Layout: layers, heads and tokens}

\section{Where the bytes go on a real accelerator}

\LedgerSection

\LabSection{a cache you can prove correct}

\HappenedSection{537 megabytes per decode step}

\FrontierSection{}

\ExercisesSection
